\section{升级为完整覆盖：先重建三条桥，再重启估值}

中国船舶案例的可复用之处是“先建立经营桥，再让估值继承模型”，不是先决定一个目标价。雅化正式 H1 披露后，必须按以下顺序重启研究。任何一条桥无法复算，报告就维持事件研究，而不是以券商 EPS、同业倍数或现价替代自有模型。

\begin{exhibitbox}[雅化完整覆盖的三条强制桥：输入、公式与失效条件]
\begin{tabularx}{\exhibitboxwidth}{L{2.2cm}X X X}
\toprule
桥 & 最低可复算式 & 正式 H1 的必要输入 & 不能满足时的处理 \\
\midrule
锂盐盈利桥 & 销量 $\times$ 实现 ASP $-$ 原料/加工成本 $=$ 毛利；再经费用、税和少数股东至归母 & H1 收入/销量、产品结构、实际 ASP 或可互证口径、原料/加工成本或分部毛利、库存及资源实际供给 & 不给资源成本优势、价格弹性或 H2 吨利润任何基准信用 \\
民爆现金桥 & 产品/服务收入 $\times$ 分部毛利率 $-$ 费用 $=$ 分部贡献；再以回款和资本占用校验现金 & 产品/爆破服务收入与毛利、订单/项目进度、龙海安防并表范围、应收和经营现金流 & 不把 2025 年高毛利率直接外推为 2026 年的稳定器 \\
利润质量桥 & 经营现金流 / 归母净利润；并追踪应收、存货、合同资产和资本开支相对收入的变化 & H1 现金流量表、资产负债表、营运资本附注和公司解释 & 不发布正常化 PE、SOTP 或基于低 PE 的行动结论 \\
\bottomrule
\end{tabularx}
\end{exhibitbox}
\sourcenote{桥的结构参照中国船舶案例的订单--交付--毛利--现金框架，并按雅化的锂盐销售型与民爆项目型业务改写。雅化的具体参数必须仅由正式 H1、后续公司披露和已归档原始资料填充。}

\section{正式 H1 发布后的模型作业卡}

完整覆盖不是把 H1 归母乘以二。下表把“何时可以开始建模”与“何时可以签发估值”分开：第一步是核对事实，第二步才是形成假设，第三步才允许估值。这样可以避免一份强预告同时被用作历史事实、全年预测和目标价锚。

\begin{exhibitbox}[从 H1 披露到完整覆盖的四步作业卡]
\begin{tabularx}{\exhibitboxwidth}{L{1.4cm}L{2.2cm}X X}
\toprule
步骤 & 交付 & 必须回答的问题 & 失败时的状态 \\
\midrule
1 & 披露核验包 & 归母/扣非、收入、毛利、OCF、应收、存货是否与预告及 Q1 相互勾稽？ & 停留在事件研究；不外推全年 \\
2 & H1 经营桥 & Q2 利润分别由销量、实现价、成本、民爆和非经常性项目贡献多少？ & 仅记录利润事件；不给 H2 持续性信用 \\
3 & H2 三情景 & 每种情景的销量、ASP、成本、民爆、费用、税、现金转换及失效条件是什么？ & 不发布 House 2026E/2027E EPS \\
4 & 双方法估值与 R1 审计 & 分部/周期正常化方法与业务匹配的可比方法是否都能复算，且现价、股本、目标、概率一致？ & 继续为“事件型观察”，不发布目标价或评级 \\
\bottomrule
\end{tabularx}
\end{exhibitbox}

步骤 3 的三情景必须是经营情景，而不是为 PE 表格填三个不同倍数。熊市情景应同时压低锂盐实现价/销量或提高成本，并检验民爆能否对冲部分波动；基准情景只能采用已被 H1 数据支持的假设；牛市情景应要求销量、价格、成本和现金转换共同达到更高门槛。情景概率必须合计 100\%，并在结构化模型中把每个参数映射到来源、置信度和失效条件。

\section{正式 H1 的六个判别字段：哪些信息真正改变价格含义}

\begin{exhibitbox}[下一份公告的读取顺序：先质量，后盈利，最后估值]
\begin{tabularx}{\exhibitboxwidth}{L{2.6cm}X X}
\toprule
字段 & 正向证据 & 为什么会改变当前结论 \\
\midrule
锂盐收入与销量 & 收入、销量与利润方向一致，且产品结构可解释 & 将“行业价格 $\beta$”转为公司级收入确认 \\
实现 ASP 与成本/毛利 & 价格上行之外，成本受控或毛利改善有可复算来源 & 判断矿--产--销协同是否值得获得持续性信用 \\
库存与原料供给 & 存货变动与销量、采购和价格敞口可以对应 & 区分备货、库存收益与周转压力 \\
民爆分部与并表 & 产品/服务毛利、项目回款和龙海安防范围有明确说明 & 判断民爆是下行缓冲还是仅有历史标签 \\
经营现金流与应收 & OCF 对利润的跟随改善，应收增速没有脱离收入 & 判断低 PE 是价值缺口还是现金质量折价 \\
非经常性、税与少数股东 & 归母改善不是由难以持续的项目主导 & 决定可进入 EPS 的利润比例 \\
\bottomrule
\end{tabularx}
\end{exhibitbox}
\sourcenote{雅化集团 2025 年年报、2026Q1 报告和 2026H1 业绩预告；中国船舶案例的 R0--R1 证据与模型门槛仅作研究流程参照。}

\section{当前行动与升级触发：把“不确定”转成明确的时间顺序}

在正式 H1 前，本报告的行动不变：\textbf{事件型观察，不能用价格回撤或外部低 PE 作为重启完整估值的理由。} 触发完整覆盖的必要条件是三条桥同时通过，而不是利润单项落在预告区间。届时新报告应交付独立的 2026E--2027E 收入、归母与 EPS，熊/基/牛经营假设与概率、两种匹配业务的估值法、现价隐含预期、可比公司对照、目标/公平价值范围及清晰的升级/降级行动；所有结论重新走 R0--R4，不继承本事件件的任何“通过”状态。

反过来，若正式 H1 仅确认利润、却不能解释现金流和营运资本，或者关键单位经济仍不披露，则最强结论仍是“Q2 弹性发生过，H2 持续性尚未证明”。这不是空泛的保守，而是对 36.3\%--56.2\% 年度盈利尚未由 H1 覆盖、以及三种外部 H2 要求差异的直接回应。

\begin{exhibitbox}[未来 R1 模型的签发底线：四类输出必须来自同一份证据与参数表]
\begin{tabularx}{\exhibitboxwidth}{L{2.7cm}X X}
\toprule
输出 & 必须与之勾稽的基础输入 & 禁止的替代做法 \\
\midrule
2026E--2027E 收入、归母与 EPS & H1 实际、锂盐量价成本、民爆收入毛利、费用、税、少数股东和加权股本 & H1 利润年化、取三家外部预测均值或仅用现货价格 \\
熊/基/牛与概率 & 每种情景的经营变量、现金转换、来源、置信度与失效条件 & 只给三个 PE 倍数，或让情景概率不加总为 100\% \\
公允价值/目标或范围 & 匹配业务模式的主方法、独立交叉方法、现价/股本/市值、市场与 Street 边界 & 以当前价、单一外部目标价或“零上行”充当估值 \\
行动与监控阈值 & 预期差、情景回报、现金质量、催化剂和可观测失效项 & 把价格波动或行业叙事单独当成基本面升级证据 \\
\bottomrule
\end{tabularx}
\end{exhibitbox}
